\scanpage{320}
\mspage{270}

la validité de $W(7)$ se ramène au cas où le carré $Y$ est le carré trivial

\diagram{\(\begin{tikzcd}[ampersand replacement=\&,column sep=tiny,row sep=tiny]
{}\&e\arrow[dr]\&{}\\e\arrow[ur]\arrow[dr]\&{}\&e\\{}\&e\arrow[ur]\&{}
\end{tikzcd},\)}

$e$ l'objet final de $\mathrm{Cat}$.

On peut maintenant conclure : \deleted{\unsure{Si}} On a

\diagram{\(\begin{tikzcd}[ampersand replacement=\&,column sep=small,row sep=small]
{}\&Z_1\arrow[dr,"{j_2}"]\&{}\\Z_0\arrow[ur,dash,"{i_1}"]\arrow[dr,"{i_2}"']\&{}\&Z_3\\{}\&Z_2\arrow[ur,"{j_1}"']\&{}
\end{tikzcd}\)}

(moyennant la partie directe de $W(7\ \mathrm{bis})$)

\[
a'\mathbin{\&}b'\Longrightarrow c'
\]

car si $Z_0,Z_1\deleted{\unsure{,Z_2}}$ sont $\in\mathrm{Asph}_W$, alors $i_1\in W$, donc $j_1\in W$, et si de plus $Z_2\in\mathrm{Asph}_W$, on a donc $Z_3\in\mathrm{Asph}_W$.

Si on a aussi la partie inverse \inserted{(``\unsure{inhabituelle|inconditionnelle}'')} de $W(7\ \mathrm{bis})$, alors

\[
b'\mathbin{\&}c'\Longrightarrow a',
\]

car si $Z_1$ et $Z_3$ sont \unsure{$W$-asph|$U$-asph}, alors $j_2\in W$, donc $i_2\in W$, et si de plus $Z_2$ est $W$ asph, on conclut $Z_0$ $W$-asph.

Enfin, supposons que $W$ soit \uline{stable par facteurs directs}, et $W(7\ \mathrm{bis})$ satisfait sous forme directe. Je dis que alors

\[
c'\mathbin{\&}a'\Longrightarrow b'
\]

\[
\text{i.e.}\qquad Z_0\text{ et } Z_3\in\mathrm{Asph}_W\Longrightarrow Z_1\text{ et }Z_2\in\mathrm{Asph}_W.
\]

Factorisant $i_1$ et $i_2$ en

\diagram{\(\begin{tikzcd}[ampersand replacement=\&,column sep=small,row sep=tiny]
Z_0\arrow[r,dash,"{i'_1}"]\&Z'_1\arrow[r,"{\alpha_1}"]\&Z_1,\\
Z_0\arrow[r,dash,"{i'_2}"]\&Z'_2\arrow[r,dash,"{\alpha_2}"]\&Z_2,
\end{tikzcd},\)}

avec $i'_1,i'_2$ des imm. ouvertes de Dwyer, et $\alpha_1,\alpha_2\in W$, on est ramené au cas où
% archive sheet number: 320
